% Elemento de volume em coordenadas esféricas
% Feito para o ÉguaTeX (eguatex.com.br). Uso livre (CC0 1.0).
%
% Compile e pronto: sai só a figura, em PDF.
% Para pôr no seu trabalho, use o PDF com \includegraphics ou copie
% o bloco tikzpicture para dentro de uma figure.
%
% Ângulo theta medido a partir do eixo z; phi no plano xy, a partir do eixo x.
% Alguns livros de cálculo trocam os nomes theta e phi.
\documentclass[tikz,border=4pt]{standalone}
\usepackage[T1]{fontenc}
\usepackage{lmodern}
\usetikzlibrary{patterns.meta}

% Para usar (esf cs:r=..., t=..., p=...) no seu documento, copie também estas linhas.
\tikzset{esf/r/.store in=\esfr, esf/t/.store in=\esft, esf/p/.store in=\esfp}
\tikzdeclarecoordinatesystem{esf}{%
  \tikzset{esf/.cd,#1}%
  \pgfpointxyz{\esfr*sin(\esft)*cos(\esfp)}{\esfr*sin(\esft)*sin(\esfp)}{\esfr*cos(\esft)}}

\begin{document}
\begin{tikzpicture}[
  x={(-0.4435cm,-0.3663cm)}, y={(0.9768cm,-0.1715cm)}, z={(0cm,1cm)},
  line cap=round, line join=round,
  eixo/.style={line width=0.75pt},
  aresta/.style={line width=0.8pt},
  guia/.style={line width=0.55pt, black!27}]

% Mude estes números e o desenho e os rótulos acompanham (ângulos em graus).
\def\ra{4.1}  \def\rb{5.45} % r e r + dr
\def\ta{41}   \def\tb{56}   % theta e theta + dtheta
\def\pa{47.5} \def\pb{70}   % phi e phi + dphi

\draw[eixo] (0,0,0) -- (5.06,0,0) node[below left] {$x$};
\draw[eixo] (0,0,0) -- (0,5.54,0) node[right]      {$y$};
\draw[eixo] (0,0,0) -- (0,0,4.93) node[above]      {$z$};

\fill[pattern={Lines[distance=3.6pt, line width=0.5pt]}, pattern color=black!25]
  plot[domain=\pa:\pb] (esf cs:r=\ra,t=90,p=\x) -- plot[domain=\pb:\pa] (esf cs:r=\rb,t=90,p=\x) -- cycle;
\draw[guia] plot[domain=\pa:\pb] (esf cs:r=\ra,t=90,p=\x);
\draw[guia] plot[domain=\pa:\pb] (esf cs:r=\rb,t=90,p=\x);
\draw[guia] (esf cs:r=\ra,t=90,p=\pa) -- (esf cs:r=\rb,t=90,p=\pa);
\draw[guia] (esf cs:r=\ra,t=90,p=\pb) -- (esf cs:r=\rb,t=90,p=\pb);

\draw[guia] (0,0,0) -- (esf cs:r=\ra,t=90,p=\pa);
\draw[guia] (0,0,0) -- (esf cs:r=\ra,t=90,p=\pb);
\foreach \r in {\ra,\rb} \foreach \p in {\pa,\pb}
  \draw[guia] plot[domain=\tb:90] (esf cs:r=\r,t=\x,p=\p);

\draw[guia] (0,0,0) -- (esf cs:r=\ra,t=\ta,p=\pb);
\draw[guia] (0,0,0) -- (esf cs:r=\ra,t=\tb,p=\pb);

\filldraw[aresta, fill=black!8]
  plot[domain=\pa:\pb] (esf cs:r=\ra,t=\tb,p=\x) -- plot[domain=\pb:\pa] (esf cs:r=\rb,t=\tb,p=\x) -- cycle;
\filldraw[aresta, fill=black!5]
  plot[domain=\ta:\tb] (esf cs:r=\ra,t=\x,p=\pa) -- plot[domain=\tb:\ta] (esf cs:r=\rb,t=\x,p=\pa) -- cycle;
\filldraw[aresta, fill=black!14]
  plot[domain=\pa:\pb] (esf cs:r=\rb,t=\ta,p=\x) -- plot[domain=\ta:\tb] (esf cs:r=\rb,t=\x,p=\pb)
  -- plot[domain=\pb:\pa] (esf cs:r=\rb,t=\tb,p=\x) -- plot[domain=\tb:\ta] (esf cs:r=\rb,t=\x,p=\pa) -- cycle;
\draw[guia] plot[domain=\pa:\pb] (esf cs:r=\ra,t=\ta,p=\x) -- (esf cs:r=\rb,t=\ta,p=\pb);
\draw[guia] plot[domain=\ta:\tb] (esf cs:r=\ra,t=\x,p=\pb);

\draw[aresta] (0,0,0) -- (esf cs:r=\ra,t=\ta,p=\pa);
\draw[aresta] (0,0,0) -- (esf cs:r=\ra,t=\tb,p=\pa);
\fill (esf cs:r=\ra,t=\ta,p=\pa) circle[radius=1.9pt];
\draw[eixo] plot[domain=0:\ta]   (esf cs:r=1.7,t=\x,p=\pa);
\draw[eixo] plot[domain=\ta:\tb] (esf cs:r=2.2,t=\x,p=\pa);
\draw[eixo] plot[domain=0:\pa]   (esf cs:r=1.65,t=90,p=\x);
\draw[eixo] plot[domain=\pa:\pb] (esf cs:r=2,t=90,p=\x);

\node at (esf cs:r=1.25,t={\ta/2},p=\pa) {$\theta$};
\node at (esf cs:r=2.85,t={(\ta+\tb)/2},p=\pa) {$\mathrm{d}\theta$};
\path (0,0,0) -- (esf cs:r=\ra,t=\ta,p=\pa) node[pos=0.7, left] {$r$};
\path (esf cs:r=\ra,t=\ta,p=\pa) -- (esf cs:r=\rb,t=\ta,p=\pa) node[midway, above left] {$\mathrm{d}r$};
\node at (esf cs:r=1.2,t=90,p={\pa/2}) {$\phi$};
\node at (esf cs:r=2.5,t=90,p={(\pa+\pb)/2}) {$\mathrm{d}\phi$};
\node at (esf cs:r=\rb,t={(\ta+\tb)/2},p={(\pa+\pb)/2}) {$\mathrm{d}V$};
\path (esf cs:r=\rb,t=\ta,p=\pb) -- (esf cs:r=\rb,t=\tb,p=\pb)
  node[midway, sloped, above] {$r\,\mathrm{d}\theta$};
\path (esf cs:r=\rb,t=\tb,p=\pa) -- (esf cs:r=\rb,t=\tb,p=\pb)
  node[midway, sloped, below=2pt] {$r\sin\theta\,\mathrm{d}\phi$};

\node[align=center] at (3.6cm,4.6cm) {\textit{Coordenadas esféricas}\\[2pt] $(r,\theta,\phi)$};
\node at (1.7cm,-3.2cm) {$\mathrm{d}V=r^2\sin\theta\,\mathrm{d}r\,\mathrm{d}\theta\,\mathrm{d}\phi$};
\end{tikzpicture}
\end{document}
